\section{Packed Metadata}
\label{sec:packed}

\subsection{An immutable snapshot with access-oriented structure}

Packed Metadata covers the metadata and access structure of selected
read-only subtrees. It is not a small standalone packet containing a few
attributes. The inspected format includes a manifest with authenticated
roots, paged indexes, directory-oriented metadata pages, object containers,
and file-data descriptors. Together these structures describe a complete
immutable lower view and how its contents can be read.

A logical read group is organized around a parent directory and a range of
raw entry names. Name order is a format-defined byte order, not a locale's
collation. Group metadata brings together directory entries, frequently
needed attributes, and locators for file data. A physical container can hold
more than one logical group, so logical grouping and physical object size
remain distinct decisions. Small contents may be inline, but metadata
co-location, inlining, compression, and data co-packing are separate factors
in our evaluation.

The current representation uses a PM11 manifest and IP06 hierarchical
indexes. Its GroupMeta raw page is bounded at 256 KiB. Other indexes support
different access paths, including inode lookup and data-frame location.
Bounding a page does not bound the entire snapshot or force all directory
entries into memory at mount time. Conversely, it does not make every
operation a constant number of GETs: a cold lookup can traverse several
pages, and a data read can require descriptor and payload access.

\begin{figure}[t]
\centering
\begin{tikzpicture}[
  b/.style={draw,rounded corners=2pt,align=center,text width=6.5cm,minimum height=0.56cm,font=\small,fill=black!3},
  a/.style={-{Latex[length=1.7mm]},semithick}
]
\node[b] (root) {Pinned manifest and authenticated roots};
\node[b,below=0.35cm of root] (index) {Paged directory-name index\\$(\text{parent},\text{raw-name interval})$};
\node[b,below=0.35cm of index] (group) {GroupMeta page\\entry $+$ common attributes $+$ data locator};
\node[b,below=0.35cm of group] (frame) {Authenticated frame descriptor\\container offset, length, content verification};
\node[b,below=0.35cm of frame] (data) {Payload Range read and verification};
\draw[a] (root) -- (index);
\draw[a] (index) -- (group);
\draw[a] (group) -- (frame);
\draw[a] (frame) -- (data);
\end{tikzpicture}

\caption{A representative lower access. A directory-name lookup descends
indexed ranges and obtains an entry, common attributes, and a data locator.
File reads use the authenticated frame descriptor and payload range.
Direct inode-based getattr and other access paths may use separate indexes;
this diagram is not a fixed GET-count claim.}
\label{fig:packed}
\end{figure}

\subsection{Lookup, attributes, enumeration, and read}

For a name lookup, the lower reader identifies the directory's relevant
range and traverses the authenticated index to its metadata page. The
returned record can supply the common attributes and location information
needed by subsequent operations. A direct getattr can instead use an inode
index carrying hot attributes; it need not fetch a whole directory group.
Figure~\ref{fig:packed} illustrates a typical name-based access, not all
possible paths.

Enumeration uses ordered ranges and page navigation to avoid requiring a
single unbounded directory object. At the snapshot-reader layer, weighted
indexes also support directory-ordinal access. These capabilities must be
distinguished from the Entity bridge's current limits: the bridge has a
bounded whole-directory interface and cannot yet expose arbitrary-sized
directory enumeration. Large-directory snapshot microbenchmarks will
therefore be reported separately from complete workspace measurements.

A file-data locator identifies an inline value or a referenced data region.
For an external frame, the reader resolves and validates its descriptor,
requests the applicable object range, and verifies the returned contents
according to the snapshot format. Failed authentication or an unavailable
referenced object is an error; it must not silently become a successful
negative name lookup. An index's absence result and a storage failure are
different outcomes.

For these supported lower queries, the needed namespace, attributes, and
data locations reside in the snapshot rather than in per-file records in
the mutable metadata database. The writable overlay may still check its
private state first. We will instrument these two kinds of database traffic
separately and test cold lower access while denying the lower-record path.
Counting only cache hits would not establish this property.

\subsection{Publication, ownership, and construction cost}

Snapshot production begins from a fixed source view. The producer groups
entries, builds metadata and data objects, and constructs the indexes and
manifest. A safe publication contract requires all referenced dependencies
to be durable and validated before the new binding becomes visible. Old
readers must retain their previous immutable view until it is safe to
reclaim. A source that changes during capture requires an explicit
consistency mechanism; an immutable output format alone does not supply it.

The repository contains producer, source-import, binding, journal, reader-pin,
and collection components. The binding proof explicitly does not, by itself,
prove source-view consistency or the closure of all uploaded dependencies.
We therefore treat end-to-end publication and failure recovery as obligations
for the selected integration, not as consequences of a successful manifest
serialization. Ordinary close or fsync does not automatically rebuild a
packed snapshot.

Publication can be expensive. Reading source metadata, sorting names,
compressing data, uploading objects, and creating authenticated indexes all
consume resources. The relevant comparison includes this cost and any
rebuild or reclamation work over the snapshot's actual reuse lifetime.
For each additive resource metric $c$, we will measure
\begin{equation}
C_c(K)=C_{c,\mathrm{build}}+C_{c,\mathrm{publish}}
       +\sum_{i=1}^{K}C_{c,i}+C_{c,\mathrm{reclaim}}.
\end{equation}
CPU-seconds, transferred bytes, and monetary cost are separate metrics;
concurrent task wall-clock times are not summed into a batch makespan.

\subsection{What the layout must earn}

Immutable metadata on object storage is already a strong existing design
space~\cite{deltafs2021,cvmfs2011,nydus_rafs}. The specific case for \sys rests
on the cost of its directory-oriented joint records and paged access under
the writable-upper contract. We will compare against an immutable sorted-KV
lower with identical record content and cache capacity, and against at least
one demand-loaded immutable-image system where its operation contract is
compatible. The sorted-KV reference is a mechanism control, not a claimed
reproduction of DeltaFS.

The layout can lose when accesses have little locality, object latency
dominates a cold index walk, or snapshots are rebuilt too frequently. These
outcomes will define its operating region. They should not be hidden by
reporting only warm directory scans.
